\documentclass[12pt]{article}
\usepackage[utf8]{inputenc}
\usepackage[T1]{fontenc}
\usepackage{hyperref}
\title{Cotoname Language Resources for K-12 Curriculum Development}
\author{Compiled by Claude Sonnet 5.5 (claude-sonnet-5-5) for Andr\'as Kornai}
\date{}
\begin{document}
\maketitle

\section{Introduction}
Cotoname (ISO 639-3 xcn; Glottocode coto1248) is an extinct language of northeastern Mexico and southern Texas. It belongs to the Comecrudan family. My searches found very little. The main documentation is Albert Gatschet's 1886 fieldwork, which recorded about 125 words from a man said to be of Cotoname descent. I found no curriculum, no community programme and no audio. I make no claim that a K-12 programme exists.

\section{Primary Educational Resources}
I found no primers, textbooks or school materials for Cotoname. Any curriculum would have to be built from the archival word lists listed below.

\section{Linguistic and Reference Materials}
\begin{itemize}
\item Wikipedia, ``Comecrudan languages'': family overview. \url{https://en.wikipedia.org/wiki/Comecrudan_languages}
\item Wikipedia, ``Comecrudo language'': the better-documented sister language. Jean Louis Berlandier recorded a list of 148 words in 1829. \url{https://en.wikipedia.org/wiki/Comecrudo_language}
\item Wikipedia, ``Pakawan languages'': discusses the proposed wider grouping, which Campbell challenged. It cites Swanton's 1940 Bureau of American Ethnology bulletin, which contains Cotoname and Comecrudo data. \url{https://en.wikipedia.org/wiki/Pakawan_languages}
\item Smithsonian National Anthropological Archives, finding aid for Gatschet's manuscript vocabularies of the Comecrudo and Cotoname dialects (MS 297). \url{https://sirismm.si.edu/EADpdfs/NAA.MS297.pdf}
\end{itemize}

\section{Religious and Community Resources}
I found no religious texts or community organisations for Cotoname in search results.

\section{Digital and Audio Resources}
I found no audio or software. Gatschet's manuscripts exist only as archival finding aids in the results.

\section{Curriculum Development Context}
The language is extinct and documentation is limited to a short word list (about 125 words). Any K-12 use would be limited to cultural and historical awareness, with the sources above cited as the documentary basis. It should not claim lexical or grammatical coverage that does not exist.
\end{document}